\backhead{EXERCISES}{khred}
\begin{exlist}
\item Observed: 20, 25, 30, 35; predicted: 22, 24, 33, 34. Find the bias, MAE and RMSE.\exans{0.75; 1.75; 1.94}
\item Compute the NSE of the same data.\exans{0.880}
\item A run has $P=640$, $I=120$, $ET=585$, $D=140$, $Q=8$, $\Delta S=22$ (mm). Find the residual and say whether the balance closes.\exans{$R=5$ mm; no, 5 mm is missing}
\item The mean residual of 25 runs is 0.8 mm with standard deviation 2.0 mm. Is it different from zero at the 5 \% level ($t_{0.975,24}=2.064$)?\exans{t = 2.0; no}
\item Find the 95 \% interval of a mean of 16 values with mean 74 and standard deviation 8 ($t_{0.975,15}=2.131$).\exans{$[69.74,\ 78.26]$}
\item Compute the bias, RMSE and NSE for observed (410, 455, 390, 500, 470, 430, 380, 460) and predicted (402, 470, 385, 487, 479, 421, 395, 449).\exans{-0.88; 11.13; 0.918}
\item If the NSE of a model is 0, what does it mean? If negative?\exans{as good as the mean; worse than the mean}
\item A slope has $b=0.82$ and standard error 0.05 with 10 degrees of freedom ($t=2.228$). Give the 95 \% interval and say whether $b=1$ is inside it.\exans{$[0.709,\ 0.931]$; no}
\item A mulch experiment shows a positive residual of +30 mm in the mulched runs and 0 in the bare runs. Where could the water have gone?\exans{uncounted evaporation from the mulch layer}
\item Why does a bootstrap interval need no assumption of normality?\exans{it uses the empirical distribution of the data by resampling}
\end{exlist}
